% Part: incompleteness
% Chapter: introduction
% Section: definitions

\documentclass[../../../include/open-logic-section]{subfiles}

\begin{document}

\olfileid{inc}{int}{def}

\olsection{Definisi}

Kanggo nindakake proyek Hilbert kang arep
ngformalkan matematika lan nuduhake yen
formalisasi iku konsisten lan jangkep,
langkah kapisan yaiku milih basa, kerangka
logis, lan sistem aksioma. Kanggo ancas kita,
upamane matematika bisa diformalkan ing basa
orde kapisan. Tegese, ana himpunan !!{constant}s,
!!{function}s, lan !!{predicate}s kang, bebarengan
karo panggandheng lan kuantor logika orde kapisan,
ngidini kita nyatakake pratelan matematika.
Akeh wong sarujuk yen basa kaya mangkono ana:
basa teori himpunan, kang mung nduwé $\in$
minangka simbol nonlogis. Sepisanan, pratelan
yen basa kang prasaja iku bisa nyatakake akèh
banget bab mesthi katon ora masuk akal; perlu
pakaryan kang akèh kanggo netepake yen meh
kabeh matematika bisa dinyatakake nganggo
kosakata kang winates iki. Supaya prasaja,
saiki kita watesi rembugan marang aritmetika,
yaiku pérangan matematika kang ngurusi
bilangan asli~$\Nat$. Basa kang lumrah kanggo
nyatakake fakta aritmetika yaiku~$\Lang L_A$.
$\Lang L_A$ ngemot siji !!{predicate} rong papan~$<$,
siji !!{constant}~$\Obj 0$, siji !!{function}
siji papan~$\prime$, lan loro !!{function}s
rong papan~$+$ lan~$\times$.

\begin{defn}
Himpunan !!{sentence}s~$\Gamma$ diarani
\emph{teori} yen katutup tumrap akibat semantis,
yaiku yen $\Gamma = \Setabs{!A}{\Gamma \Entails
!A}$.
\end{defn}

Ana rong cara gampang kanggo netepake teori.
Cara kapisan yaiku minangka himpunan !!{sentence}s
kang bener ing sawijining !!{structure}. Tuladhane,
gatekna !!{structure} kanggo $\Lang L_A$ kang
!!{domain}e yaiku~$\Nat$ lan kabeh simbol
nonlogis diinterpretasi kaya kang lumrah.

\begin{defn}\ollabel{def:standard-model}
\emph{Model baku aritmetika} yaiku
!!{structure}~$\Struct N$ kang ditegesi kaya mangkene:
\begin{enumerate}
\item $\Domain N = \Nat$
\item $\Assign{\Obj 0}{N} = 0$
\item $\Assign{\Obj \prime}{N}(n) = n + 1$ kanggo kabeh $n \in \Nat$
\item $\Assign{\Obj +}{N}(n, m) = n + m$ kanggo kabeh $n, m \in \Nat$
\item $\Assign{\Obj \times}{N}(n, m) = n\cdot m$ kanggo kabeh $n, m \in \Nat$
\item $\Assign{\Obj <}{N} = \Setabs{\tuple{n, m}}{n \in \Nat, m \in
  \Nat, n < m}$
\end{enumerate}
\end{defn}

Gatekna bédane `$\times$' lan `$\cdot$':
$\times$ iku simbol ing basa aritmetika.
Mesthi wae, simbol iku dipilih supaya
ngelingake kita marang perkalian, nanging
$\times$ dudu operasi perkalian; iku simbol
fungsi rong papan (sacara resmi, $\Obj f^2_1$).
Déné $\cdot$~\emph{pancen} fungsi perkalian lumrah.
Yen ana tulisan kaya $n \cdot m$, tegesé asil
perkalian bilangan $n$ lan~$m$; yen ana tulisan
kaya $x \times y$, kang dirembug yaiku suku
ing basa aritmetika. Ing model baku, simbol
fungsi~$\times$ diinterpretasi minangka fungsi
$\cdot$ ing bilangan asli. Kanggo panjumlahan,
kita nggunakake $+$ minangka simbol fungsi
ing basa aritmetika lan uga minangka fungsi
panjumlahan ing bilangan asli. Ing kene konteks
kang nemtokake tegesé.

\begin{defn}
Teori \emph{aritmetika bener} yaiku himpunan
!!{sentence}s kang dicukupi ing model baku
aritmetika, yaiku
\[
\Th{TA} = \Setabs{!A}{\Sat{N}{!A}}.
\]
\end{defn}

$\Th{TA}$ iku teori, amarga yen
$\Th{TA} \Entails !A$, $!A$ dicukupi ing
saben !!{structure} kang nyukupi~$\Th{TA}$.
Amarga $\Sat{N}{\Th{TA}}$, mula
$\Sat{N}{!A}$, dadi $!A \in \Th{TA}$.

Cara liyane kanggo netepake teori~$\Gamma$
yaiku minangka himpunan !!{sentence}s kang
dadi akibat semantis saka himpunan ukara~$\Gamma_0$.
Ing kahanan iki, $\Gamma$ iku ``panutupan''
$\Gamma_0$ tumrap akibat semantis. Netepake
teori kanthi cara iki mung migunani yen
$\Gamma_0$ ditemtokake kanthi cetha, umpamane
!!{element}s saka~$\Gamma_0$ didhaptar.
Paling ora, $\Gamma_0$ kudu bisa diputusake,
yaiku kudu ana tes komputabel kanggo mutusake
apa !!a{sentence} kalebu unsur~$\Gamma_0$
utawa ora. !!{sentence}s ing~$\Gamma_0$
diarani \emph{aksioma} tumrap~$\Gamma$,
dene $\Gamma$ diarani \emph{diaksiomatisasi}
dening~$\Gamma_0$.

\begin{defn}
Teori~$\Gamma$ \emph{diaksiomatisasi}
dening~$\Gamma_0$ yen lan mung yen
\[
\Gamma = \Setabs{!A}{\Gamma_0 \Entails !A}
\]
\end{defn}

\begin{defn}
Teori $\Th{Q}$ kang diaksiomatisasi dening
ukara-ukara ing ngisor iki dikenal minangka
``$\Th{Q}$ Robinson'' lan dadi teori
aritmetika kang prasaja banget.
\begin{align*}
& \lforall[x][\lforall[y][(\eq[x'][y'] \lif \eq[x][y])]] \tag{$!Q_1$}\\
& \lforall[x][\eq/[\Obj 0][x']] \tag{$!Q_2$}\\
& \lforall[x][(\eq[x][\Obj 0] \lor \lexists[y][\eq[x][y']])] \tag{$!Q_3$}\\
& \lforall[x][\eq[(x + \Obj 0)][x]] \tag{$!Q_4$}\\
& \lforall[x][\lforall[y][\eq[(x + y')][(x + y)']]] \tag{$!Q_5$}\\
& \lforall[x][\eq[(x \times \Obj 0)][\Obj 0]] \tag{$!Q_6$}\\
& \lforall[x][\lforall[y][\eq[(x \times y')][((x \times y) + x)]]] \tag{$!Q_7$}\\
& \lforall[x][\lforall[y][(x < y \liff \lexists[z][\eq[(z' + x)][y]])]] \tag{$!Q_8$}
\end{align*}
Himpunan !!{sentence}s $\{!Q_1, \dots, !Q_8\}$
iku aksioma $\Th{Q}$, mula $\Th{Q}$ ngemot
kabeh !!{sentence}s kang dadi akibat semantise:
\[
\Th{Q} = \Setabs{!A}{\{!Q_1, \dots, !Q_8\} \Entails !A}.
\]
\end{defn}

\begin{defn}
Upamane $!A(x)$ iku !!a{formula} ing $\Lang L_A$
kanthi variabel bebas~$x$ lan $y_1$, \dots, $y_n$.
Mula saben !!{sentence} kanthi wujud
\[
\lforall[y_1][\dots\lforall[y_n][((!A(\Obj 0) \land \lforall[x][(!A(x)
\lif !A(x'))]) \lif \lforall[x][!A(x)])]]
\]
iku instansiasi saka \emph{skema induksi}.

\emph{Aritmetika Peano}~$\Th{PA}$ yaiku teori
kang diaksiomatisasi dening aksioma $\Th{Q}$
bebarengan karo kabeh instansiasi skema induksi.
\end{defn}

\begin{explain}
Saben instansiasi skema induksi bener
ing~$\Struct{N}$. Iki paling gampang dideleng
yen !!{formula}~$!A$ mung nduwé siji
!!{variable} bebas~$x$. Banjur $!A(x)$
netepake subhimpunan~$X_{!A}$ saka~$\Nat$
ing~$\Struct{N}$. $X_{!A}$ iku himpunan
kabeh~$n \in \Nat$ saengga
$\Sat{N}{!A(x)}[s]$ nalika $s(x) = n$.
Instansiasi skema induksi kang cocog yaiku
\[
((!A(\Obj 0) \land \lforall[x][(!A(x) \lif !A(x'))]) \lif 
  \lforall[x][!A(x)]).
\]
Yen antesedene bener ing~$\Struct{N}$,
mula $0 \in X_{!A}$ lan, saben
$n \in X_{!A}$, uga $n+1$ kalebu.
Amarga $0 \in X_{!A}$, kita entuk
$1 \in X_{!A}$. Saka $1 \in X_{!A}$
kita entuk $2 \in X_{!A}$, lan sateruse.
Dadi kanggo saben $n \in \Nat$,
$n \in X_{!A}$. Iki ateges
$\lforall[x][!A(x)]$ dicukupi ing~$\Struct{N}$.
\end{explain}

$\Th{Q}$ lan $\Th{PA}$ loro-lorone teori
kang diaksiomatisasi. Pitakon gedhene,
sepira kuwat teori-teori iku? Umpamane,
apa $\Th{PA}$ bisa mbuktekake kabeh
kabeneran ngenani~$\Nat$ kang bisa
dinyatakake ing~$\Lang L_A$? Luwih
cetha maneh, apa aksioma $\Th{PA}$
bisa mutusake kabeh pitakon kang bisa
dirumusake ing~$\Lang L_A$?

Pitakon iki uga bisa diucapake mangkene:
Apa $\Th{PA} = \Th{TA}$? Mesthi wae
$\Th{TA}$ bisa mbuktekake (yaiku ngemot)
kabeh kabeneran ngenani~$\Nat$, lan bisa
mutusake kabeh pitakon kang bisa dirumusake
ing~$\Lang L_A$. Sebabe, yen $!A$ iku
!!a{sentence} ing $\Lang L_A$, mesthi
$\Sat{N}{!A}$ utawa $\Sat{N}{\lnot !A}$,
mula $\Th{TA} \Entails !A$ utawa
$\Th{TA} \Entails \lnot !A$. Teori
kaya mangkono diarani \emph{!!{complete}}.

\begin{defn}
Teori $\Gamma$ diarani \emph{!!{complete}}
yen lan mung yen kanggo saben
!!{sentence}~$!A$ ing basane,
$\Gamma \Entails !A$ utawa
$\Gamma \Entails \lnot !A$.
\end{defn}

\begin{explain}
Miturut Teorema Kelengkapan,
$\Gamma \Entails !A$ yen lan mung yen
$\Gamma \Proves !A$. Dadi $\Gamma$
jangkep yen lan mung yen kanggo saben
!!{sentence}~$!A$ ing basane,
$\Gamma \Proves !A$ utawa
$\Gamma \Proves \lnot !A$.
\end{explain}

Pitakon liyane yaiku: Apa ana prosedur komputasional
kanggo mriksa apa !!a{sentence} kalebu~$\Th{TA}$,
$\Th{PA}$, utawa malah mung~$\Th{Q}$?
Iki bisa digawe luwih cetha kanthi netepake
kapan sawijining himpunan (umpamane himpunan
!!{sentence}s) iku !!{decidable}.

\begin{defn}
Himpunan $X$ diarani \emph{!!{decidable}}
yen lan mung yen ana prosedur komputasional
kang, kanggo input~$x$, ngasilake~$1$
yen $x \in X$ lan $0$ yen ora.
\end{defn}

Dadi pitakone: Apa $\Th{TA}$
($\Th{PA}$, $\Th{Q}$) iku !!{decidable}?

Wangsulan kanggo kabeh pitakon iki yaiku: ora.
Ora siji wae saka teori-teori iki kang bisa
diputusake. Nanging kedadeyan iki ora mung
kanggo teori-teori iki. Sejatine, \emph{sembarang}
teori kang nyukupi syarat tartamtu kena asil
kang padha. Salah sawijining syarat iku,
kang dicukupi dening $\Th{Q}$ lan $\Th{PA}$,
yaiku teori kasebut diaksiomatisasi dening
himpunan aksioma kang !!a{decidable}.

\begin{defn}
Teori diarani \emph{!!{axiomatizable}}
yen diaksiomatisasi dening himpunan
aksioma kang !!a{decidable}.
\end{defn}

\begin{ex}
Saben teori kang diaksiomatisasi dening himpunan
winates !!{sentence}s iku !!{axiomatizable},
amarga saben himpunan winates iku !!{decidable}.
Mula, tuladhane, $\Th{Q}$ iku !!{axiomatizable}.

Teori kang diaksiomatisasi nganggo skema kaya~$\Th{PA}$
uga !!{axiomatizable}. Kanggo mriksa apa $!B$
kalebu aksioma~$\Th{PA}$, yaiku ngitung
fungsi $\Char{X}$ kang nduwé
$\Char{X}(!B) = 1$ yen $!B$ aksioma~$\Th{PA}$
lan $= 0$ yen ora, kita bisa nindakake iki:
Kapisan, priksa apa $!B$ salah siji aksioma~$\Th{Q}$.
Yen iya, wangsulane ``ya'' lan nilaine
$\Char{X}(!B) = 1$. Yen ora, priksa apa iku
instansiasi skema induksi. Iki bisa ditindakake
kanthi sistematis; ing kene cara kang paling
gampang dideleng kaya mangkene: Saben instansiasi
skema induksi diwiwiti sawetara kuantor universal,
banjur sub-!!{formula} kang wujude implikasi.
Konsekuen implikasi iku
$\lforall[x][!A(x, y_1, \dots, y_n)]$, ing ngendi
$x$ lan $y_1$, \dots, $y_n$ iku kabeh variabel
bebas saka~$!A$, dene kuantor-kuantor wiwitan
ing~$!B$ ngiket variabel~$y_1$, \dots,~$y_n$.
Sawisé ngangkat~$!A$ iki lan mriksa yen variabel
bebase cocog karo variabel kang diiket kuantor
universal ing ngarep lan~$\lforall[x]$,
kita mriksa apa anteseden implikasi iku cocog karo
\[
!A(\Obj 0, y_1, \dots, y_n) \land \lforall[x][(!A(x, y_1, \dots, y_n)
\lif !A(x', y_1, \dots, y_n))]
\]
Yen cocog, $!B$ iku instansiasi skema induksi;
yen ora cocog, $!B$ dudu instansiasi iku.
\end{ex}

Kanggo mangsuli pitakon iki---lan pitakon
kang luwih umum ngenani teori endi kang jangkep
utawa bisa diputusake---definisi ing ngisor iki
uga migunani. Elinga yen himpunan $X$ iku
!!{enumerable} yen lan mung yen kosong utawa
ana fungsi~$f \colon \Nat \to X$ kang
!!a{surjective}. Fungsi kaya mangkono diarani
enumerasi saka~$X$.

\begin{defn}
Himpunan $X$ diarani \emph{!!{computably enumerable}}
(disingkat !!{c.e.}) yen lan mung yen kosong
utawa nduwé enumerasi komputabel.
\end{defn}

Saliyane !!{axiomatizability}, syarat liyane
kanggo teori kang kena teorema-teorema ora
jangkep yaiku cukup kuwat kanggo mbuktekake
fakta dhasar ngenani fungsi komputabel lan
relasi !!{decidable}. Kang diarani ``fakta
dhasar'' yaiku !!{sentence}s kang nyatakake
nilai fungsi komputabel kanggo saben argumene.
Kang diarani ``cukup kuwat'' yaiku teori kasebut
ngemot ukara-ukara iki minangka teoremane.
Tuladhane, gatekna fungsi komputabel kang lumrah:
panjumlahan. Nilai $+$ kanggo argumen $2$
lan $3$ yaiku~$5$, yaiku $2+3 = 5$.
Ukara ing basa aritmetika kang nyatakake yen
nilai $+$ kanggo argumen $2$ lan $3$ iku~$5$
yaiku: $(\num{2} + \num{3}) = \num{5}$.
Lan, umpamane, $\Th{Q}$ mbuktekake ukara iki.
Luwih umum, kanggo saben fungsi komputabel
$f(x_1, x_2)$, kita kepengin ana
!!a{formula}~$!A_f(x_1, x_2, y)$ ing~$\Lang{L_A}$
saengga $\Th{Q} \Proves !A_f(\num{n_1},
\num{n_2}, \num{m})$ yen $f(n_1, n_2) = m$.
Kanthi mangkono, $\Th{Q}$ mbuktekake yen
nilai~$f$ kanggo argumen $n_1$, $n_2$
iku~$m$. Malah kita mbutuhake luwih saka iku:
teori iku uga kudu mbuktekake yen ora ana
bilangan liya kang dadi nilai~$f$ kanggo
argumen $n_1$,~$n_2$. Bab kang padha uga
kanggo relasi !!{decidable}. Iki bakal
ditegesi kanthi cetha ing rong definisi
ing ngisor iki.

\begin{defn}
!!{formula}~$!A(x_1, \dots, x_k, y)$
\emph{!!{represents}} fungsi
$f\colon \Nat^k \to \Nat$ ing~$\Gamma$
yen lan mung yen saben $f(n_1, \dots, n_k) = m$,
mangka
\begin{enumerate}
\item $\Gamma \Proves !A(\num{n_1}, \dots, \num{n_k}, \num{m})$, lan
\item $\Gamma \Proves \lforall[y](!A(\num{n_1}, \dots, \num{n_k},
y) \lif y = \num{m})$.
\end{enumerate}
\end{defn}

\begin{defn}
% OLPL-167: Definisi sumber kanggo representasi relasi nggunakake Gamma ing klausa bukti nanging ora ngenalake parameter "ing Gamma" kaya ing definisi fungsi sadurunge.
!!{formula}~$!A(x_1, \dots, x_k)$
\emph{!!{represents}} relasi
$R \subseteq \Nat^k$ yen lan mung yen,
\begin{enumerate}
\item saben $R(n_1, \dots, n_k)$ bener,
  $\Gamma \Proves !A(\num{n_1},
\dots, \num{n_k})$, lan
\item saben $R(n_1, \dots, n_k)$ ora bener,
  $\Gamma \Proves \lnot
!A(\num{n_1}, \dots, \num{n_k})$.
\end{enumerate}
\end{defn}

Teori diarani ``cukup kuwat'' supaya
teorema-teorema ora jangkep bisa ditrapake
yen teori iku !!{represents} kabeh fungsi
komputabel lan kabeh relasi !!{decidable}.
$\Th{Q}$~lan perluasane nyukupi syarat iki,
nanging kita isih perlu wektu kanggo netepake
iku---iki fakta kang ora prasaja ngenani apa
kang bisa dibuktekake dening $\Th{Q}$,
lan angel dituduhake amarga $\Th{Q}$ mung
nduwé sithik aksioma minangka landhesan
mbuktekake kabeh fakta kasebut. Nanging
$\Th{Q}$ iku teori kang ringkih banget.
Mula sanadyan angel mbuktekake yen
$\Th{Q}$ nggambarake kabeh fungsi komputabel,
akeh teori kang narik kawigaten luwih
kuwat tinimbang~$\Th{Q}$, yaiku bisa
mbuktekake luwih akeh tinimbang $\Th{Q}$.
Yen $\Th{Q}$ mbuktekake sawijining bab,
teori kang luwih kuwat uga bisa; amarga
$\Th{Q}$ nggambarake kabeh fungsi komputabel,
saben teori kang luwih kuwat uga bisa.
Iki ateges akeh teori kang narik kawigaten
nyukupi syarat teorema-teorema ora jangkep.
Dadi upaya kita kang angel iki bakal ana
asilé, amarga nuduhake yen teorema-teorema
ora jangkep ditrapake marang akèh teori.
Mesthi wae, saben teori kang arep
ngformalkan ``kabeh matematika'' kudu
mbuktekake kabeh kang dibuktekake
dening $\Th{Q}$, amarga paling ora teori
iku kudu bisa nyakup asil komputasi dhasar.
Mula saben teori kang bisa dadi calon
teori kanggo ``kabeh matematika'' uga
kena teorema-teorema ora jangkep.


\end{document}
